\section{Cote 134-8: the linear lemmas, normalization, and a proved range of the proposed theorem}

\subsection{The recovered question and the scope of its predicates}
The primary source is Grothendieck's \emph{Pursuing stacks}, cote
134-8, manuscript footers 47--54, in the official Montpellier PDF
\url{https://grothendieck.umontpellier.fr/134-8.pdf}.  This PDF has a
cover: its physical page number is the manuscript footer plus one.
The displayed \emph{Thm(?)} is on footer 48, physical page 49.
The original uses a connective chain-complex target, denoted with a
subscript on \(D\); we do not replace it by the unbounded derived
category.  The statements proved below use the unambiguous modern
notation \(D_{\geq0}(k)\), with homological differentials of degree
\(-1\).

On footer 47 the scalar functor \(\kappa(U)=\mathbf1\otimes_k U\)
is required to commute with inverse limits and is given a left adjoint
\(\varphi\).  The separately mentioned right adjoint of \(\kappa\)
is \(\operatorname{Hom}_{\mathcal M}(\mathbf1,-)\).  The setup treats
\(\kappa\) as exact.  There are enough projectives, the canonical
map \(L\varphi(\mathbf1)\to k\) is an isomorphism, and a monomorphism
\(\alpha:\mathbf1\to J\) is sought with
\[
 X\longrightarrow J\otimes X\text{ monic},\qquad
 L\varphi(J\otimes X)=0\quad(X\in\mathcal M).
 \tag{134.1}
\]
The vanishing includes degree zero.  The proposed conclusion is that
localizing \(\mathcal M\) at the maps inverted by \(L\varphi\)
gives the connective derived category of scalar modules.  Footer 48
then adds projective generators of finite presentation and sketches a
chain-complex functor and its right adjoint.  Its right margin also
records a projectivity condition on the scalar modules of morphisms.
We state this extra condition explicitly in BR-13405.

The supplied DER\_LOCALIZERS intake reports results C.1--C.29,
including a counterexample outside the finitely presented-generator
case.  Their original proofs are absent.  The present proofs therefore
establish the claims they state directly; they do not certify the
unseen counterexample or silently identify its hypotheses with these
handwritten predicates.  No RH or GRH input is needed for this section.

\subsection{BR-13401. The weakened form of Lemme 1 is valid}
Let \(\mathcal M,\mathcal A\) be abelian categories, with enough
projectives in \(\mathcal M\), and let
\(\Lambda:\mathcal M\to\operatorname{Ch}_{\geq0}(\mathcal A)\)
be an exact additive functor.  Put \(\varphi=H_0\Lambda\).
Assume that, for every projective \(P\in\mathcal M\),
\[
 H_j(\Lambda P)=0\quad(j>0).
 \tag{134.2}
\]
Then \(\varphi\) is right exact and there is a natural isomorphism
\[
 L\varphi(X)\simeq\Lambda X
 \quad\text{in }D_{\geq0}(\mathcal A),\qquad X\in\mathcal M.
 \tag{134.3}
\]
It is unnecessary to require \(\Lambda P\) to be a projective object
of the chain-complex category.  This is the weakened condition written
beside Lemme 1 on manuscript footer 53, physical page 54.

\paragraph{Complete proof.}
Exactness of \(\Lambda\) and right exactness of \(H_0\) imply right
exactness of \(\varphi\).  Choose a projective resolution
\(P_\bullet\to X\) in nonnegative homological degrees.  Form the
first-quadrant double complex
\(B_{a,b}=(\Lambda P_a)_b\), with the usual sign in one
differential so that the differentials anticommute.  Its total complex
uses only the finite direct sum \(\bigoplus_{a+b=n}B_{a,b}\) in each
degree \(n\).

The augmentations \(P_\bullet\to X\) give a map
\(\operatorname{Tot}B\to\Lambda X\).  In every fixed internal degree
\(b\), its augmented horizontal complex is exact, since the degree
\(b\) functor of \(\Lambda\) is exact.  The first-quadrant filtration
therefore proves this map is a quasi-isomorphism.  In the other
direction, the natural augmentations
\(\Lambda P_a\to H_0(\Lambda P_a)[0]\) give a map
\(\operatorname{Tot}B\to\varphi(P_\bullet)\).  Its vertical homology
is an isomorphism in degree zero and zero in positive degree by
(134.2).  The other first-quadrant filtration proves that it too is a
quasi-isomorphism.  These assertions can equivalently be checked by
successive filtration quotients in each total degree; no convergence
of an infinite diagonal is being assumed.  The complex
\(\varphi(P_\bullet)\) computes \(L\varphi(X)\), giving (134.3).
Comparison maps between projective resolutions are unique up to chain
homotopy; applying the additive functor and totalization preserves
these homotopies.  The resulting isomorphism is consequently natural
and independent of the resolution. \(\square\)

\subsection{BR-13402. Projective objects of the connective chain category}
Suppose \(\mathcal A\) is abelian with enough projectives and
coproducts, and coproducts of projectives are projective.  For
\(n\geq1\), let \(D^n(Q)\) be the complex with \(Q\) in degrees
\(n,n-1\), identity differential, and zero elsewhere.  A complex
\(P\) is projective as an object of
\(\operatorname{Ch}_{\geq0}(\mathcal A)\) if and only if it is
isomorphic to
\[
 U[0]\oplus\bigoplus_{n\geq1}D^n(Q_{n-1}),
 \tag{134.4}
\]
where \(U,Q_j\) are projective in \(\mathcal A\).  The sum here is
locally finite in the chain degree.  In particular, a projective
complex is not the same as an arbitrary complex with projective terms.

\paragraph{Complete proof.}
There are natural identifications
\[
 \operatorname{Hom}_{\mathrm{Ch}}(Q[0],C)
   =\operatorname{Hom}_{\mathcal A}(Q,C_0),\qquad
 \operatorname{Hom}_{\mathrm{Ch}}(D^nQ,C)
   =\operatorname{Hom}_{\mathcal A}(Q,C_n).
\]
Consequently the indicated building blocks are projective.  Their
coproducts are projective, proving one implication.

For the converse, choose projective epimorphisms \(Q_n\to P_n\)
for all \(n\geq0\).  The preceding identifications give a degreewise
epimorphism
\[
 Q_0[0]\oplus\bigoplus_{n\geq1}D^nQ_n\longrightarrow P.
\]
It splits because \(P\) is projective.  Thus the terms of \(P\)
are projective, \(H_jP=0\) for \(j>0\), and \(H_0P\) is
projective, all by taking the corresponding direct summands of the
source.  Write \(B_j=\operatorname{im}(d:P_{j+1}\to P_j)\).
The sequence \(0\to B_0\to P_0\to H_0P\to0\) splits, so
\(B_0\) is projective.  Inductively,
\(0\to B_j\to P_j\to B_{j-1}\to0\) splits for every \(j\geq1\),
because \(H_jP=0\) and \(B_{j-1}\) is projective.  These splittings
identify \(P_0\) with \(H_0P\oplus B_0\) and \(P_j\) with
\(B_j\oplus B_{j-1}\), with differential the identity on the second
factor and zero on the first.  This is (134.4), with \(U=H_0P\).
\(\square\)

\subsection{BR-13403. Unit normalization and the exact remaining counit test}
Let \(k\) be a commutative ring.  Assume the following precise data:
\(\mathcal M\) is a cocomplete \(k\)-linear abelian category with
enough projectives; \(\kappa(U)=\mathbf1\otimes_k U\) is an exact
scalar functor with left adjoint \(\varphi\); and there is an exact
\(k\)-linear adjunction
\[
 \Lambda:\mathcal M\rightleftarrows
 \operatorname{Ch}_{\geq0}(k):G
 \tag{134.5}
\]
with both functors exact, \(H_0\Lambda=\varphi\), and (134.2).
Assume also that the canonical counit
\(\Lambda G(k[0])\to k[0]\) is a quasi-isomorphism.  Equivalently,
under the identification below, this is the canonical normalization
\(L\varphi(\mathbf1)\to k\), including degree zero.

Set \(W=\Lambda^{-1}(\mathrm{qis})\).  Then (134.5) induces
mutually inverse equivalences on the localizations, with the induced
unit and counit, if and only if
\[
 L\varphi(GP)=0
 \quad\text{for every acyclic projective object }
 P\in\operatorname{Ch}_{\geq0}(k).
 \tag{134.6}
\]
Thus (134.6) is an exact residual test for this specified adjunction,
not merely a sufficient list of examples.  It is the question left at
the bottom of manuscript footer 54, physical page 55.

\paragraph{Complete proof.}
For every scalar module \(U\) and every \(X\in\mathcal M\),
the adjunctions give natural isomorphisms
\[
 \operatorname{Hom}_{\mathcal M}(X,G(U[0]))
 =\operatorname{Hom}_{\mathrm{Ch}}(\Lambda X,U[0])
 =\operatorname{Hom}_k(\varphi X,U)
 =\operatorname{Hom}_{\mathcal M}(X,\kappa U).
\]
Yoneda identifies \(G(U[0])\) with \(\kappa U\), and identifies
the two counits.  If \(U\) is projective, it is a summand of a free
module.  Both \(\kappa\) and \(\Lambda\) preserve coproducts:
the first by its scalar-tensor construction, the second because it is
a left adjoint.  Therefore the normalization for \(k\) implies that
\(\Lambda G(U[0])\to U[0]\) is a quasi-isomorphism for every
projective \(U\), by coproducts and retracts.  The coproduct of these
quasi-isomorphisms is a quasi-isomorphism in modules, since direct
sums are exact.

Every projective chain object splits by BR-13402 as \(U[0]\oplus P\)
with \(U\) projective and \(P\) acyclic projective.  Additivity of
both adjoints means that its counit is a quasi-isomorphism exactly when
the counit on \(P\) is one.  Since \(P\) is acyclic, the latter
condition is precisely \(\Lambda GP\simeq0\), equivalently (134.6)
by BR-13401.  The exact-adjunction criterion F-SM17 proves that counits
on all projective chain objects are quasi-isomorphisms if and only if
all counits are, and that this is equivalent to the stated localized
adjunction equivalence.  This proves both implications. \(\square\)

\subsection{BR-13404. The proposed equivalence holds over fields and over \(\mathbf Z\)}
Retain all the data of BR-13403.  Assume \(\mathcal M\) has enough
injectives and the tensor interval (134.1) exists.  Suppose there is
an integer \(r\geq0\) such that every projective \(k\)-module has
injective dimension at most \(r\).  Then (134.6) holds and the exact
adjunction induces an equivalence
\[
 W^{-1}\mathcal M\simeq D_{\geq0}(k).
 \tag{134.7}
\]
In particular this conclusion holds for a field \(k\) and for
\(k=\mathbf Z\), once the specified exact adjunction has been
constructed.  These are the rational and integral scalar rings used
in the programme.  The result does not assert an unbounded derived
equivalence or eliminate the construction hypotheses of BR-13403.

\paragraph{Complete proof.}
If \(I\) is injective in \(\mathcal M\), its monomorphism
\(I\to J\otimes I\) splits.  Applying \(L\varphi\) and (134.1)
shows \(L\varphi(I)=0\), since it is a retract of zero.  Any object
\(X\) having a finite injective resolution consequently satisfies
\(L\varphi(X)=0\).  Indeed, in a short exact sequence
\(0\to X\to I\to X'\to0\) with \(I\) injective, the triangle for
the left derived functor gives
\(L\varphi(X')\simeq L\varphi(X)[1]\).  Starting at the final
injective term of a finite resolution and working backwards proves
the assertion, including degree-zero vanishing.

Let \(P\) be an acyclic projective chain object.  BR-13402 gives
\(P=\bigoplus_{n\geq1}D^n(Q_{n-1})\) for projective scalar
modules \(Q_j\).  Choose length-\(r\) injective resolutions of the
\(Q_j\).  Apply each disk functor and take the locally finite sum
over \(j\).  In each chain degree only two of these resolutions
occur, so the resulting sequence is exact; it is a length-\(r\)
resolution of \(P\).

Each of its terms is injective as a chain object.  To see this without
an unsupported assertion about infinite sums of injectives, note that
\(D^{j+1}(I)\) is right adjoint, in the appropriate variable, to
the exact evaluation \(C\mapsto C_j\):
\[
 \operatorname{Hom}_{\mathrm{Ch}}(C,D^{j+1}I)
    =\operatorname{Hom}_k(C_j,I).
\]
It is therefore injective if \(I\) is injective.  The locally finite
sum of these disks equals their product degree by degree.  Products
of injective objects are injective, since maps into the product are
products of maps and products of surjections of abelian groups are
surjective.  This proves the asserted injectivity.

Because \(\Lambda\) is exact, its right adjoint \(G\) sends
injectives to injectives.  Because \(G\) is exact, applying it to
the length-\(r\) resolution gives a finite injective resolution of
\(GP\) in \(\mathcal M\).  The first paragraph now gives
\(L\varphi(GP)=0\).  BR-13403 proves (134.7).

Over a field, every module is injective, so \(r=0\) works.  Over
\(\mathbf Z\), every free module \(F\) has
\(0\to F\to F\otimes\mathbf Q\to(F\otimes\mathbf Q)/F\to0\),
whose two latter groups are divisible and hence injective.  Every
projective abelian group is a summand of a free one; taking summands
shows injective dimension at most one.  For completeness, the
divisibility assertion uses the elementary injectivity criterion:
a homomorphism \(n\mathbf Z\to D\) extends to \(\mathbf Z\to D\)
exactly when the prescribed image of \(n\) is divisible by \(n\).
The usual ideal criterion for injectivity then applies.  Thus \(r=1\)
works over \(\mathbf Z\). \(\square\)

\subsection{BR-13405. An explicit construction of the exact adjunction in the corollary's locally projective case}
Let \(R\) be a small \(k\)-linear category such that every
\(R(x,y)\) is a projective \(k\)-module, and let
\(\mathcal M=\operatorname{Fun}_k(R^{\mathrm{op}},\operatorname{Mod}k)\).
Assume \(\kappa(U)=T\otimes_k U\), for a fixed right \(R\)-module
\(T\), is exact and has a left adjoint \(\varphi\).  Then there
is an exact \(k\)-linear functor \(\Lambda\) representing
\(L\varphi\) on objects and an exact right adjoint \(G\), as in
(134.5).  This supplies a fully specified construction when the
finite-presentation projective generators and the additional scalar
projectivity condition on footer 48 are interpreted by such a ringoid.

\paragraph{Complete proof.}
Write \(P_x=R(-,x)\) for the right representable and set
\(D(x)=\varphi(P_x)\).  The dependence on \(x\) makes \(D\) a
left \(R\)-module.  Since \(\kappa\) is exact, its left adjoint
\(\varphi\) sends projectives to projectives, so each \(D(x)\)
is a projective scalar module.  In fact it is finitely generated:
\[
 \operatorname{Hom}_k(D(x),U)
   =\operatorname{Hom}_{\mathcal M}(P_x,\kappa U)
   =T(x)\otimes_k U
\]
commutes with arbitrary coproducts.  A projective module with this
property is finitely generated.  To verify this last assertion, write
it as a summand of \(k^{(I)}\).  Its inclusion in that direct sum
has, by the coproduct property, only finitely many nonzero coordinate
maps, so it is a summand of a finite free module.

There is a natural isomorphism \(\varphi(X)=D\otimes_R X\).
Both sides preserve coproducts and cokernels, and they agree naturally
on each representable.  Every right module has a presentation by
coproducts of representables; applying the two functors to such a
presentation proves the isomorphism and its naturality.

Take a projective resolution \(B_\bullet\to D\) by left
\(R\)-modules, choosing each \(B_a\) to be a coproduct of left
representables \(R(x,-)\).  Put
\[
 \Lambda(X)=B_\bullet\otimes_R X.
 \tag{134.8}
\]
Every degree of this functor is a sum of evaluation functors, and is
therefore exact; hence \(\Lambda\) is exact.  Its zeroth homology
is \(D\otimes_R X=\varphi X\).  For \(X=P_x\), the complex is
\(B_\bullet(x)\), whose positive homology vanishes and whose zeroth
homology is \(D(x)\).  Coproducts and retracts give the same
positive-homology vanishing for every projective \(X\).
BR-13401 proves that (134.8) represents \(L\varphi\).

Define a right \(R\)-module by
\[
 (GC)(x)=\operatorname{Hom}_{\operatorname{Ch}_{\geq0}(k)}
                    (B_\bullet(x),C).
 \tag{134.9}
\]
The ordinary tensor--Hom adjunction, degree by degree with the chain
map compatibility, gives \(\Lambda\dashv G\).  Every term
\(B_a(x)\) is a coproduct of projective scalar modules \(R(y,x)\),
and is projective.  The complex \(B_\bullet(x)\) resolves the
projective module \(D(x)\).  Its successive boundary sequences
therefore split, exactly as in BR-13402; it is a projective object of
the chain category.  It follows that (134.9) is an exact functor of
\(C\), for every \(x\).  Exactness in a functor category is
objectwise, so \(G\) is exact.  This completes the construction.
\(\square\)

\subsection{BR-13406. Positive-degree interval acyclicity is insufficient}
The degree-zero clause in (134.1) cannot be removed.  More precisely,
replace its vanishing by
\(L_i\varphi(J\otimes X)=0\) only for \(i>0\), keeping the scalar
unit normalization and all the exact-adjunction data.  The proposed
canonical equivalence need no longer hold.

\paragraph{Complete counterexample.}
Take a nonzero field \(k\), \(\mathcal M=\operatorname{Mod}k\)
with its ordinary tensor product, \(\kappa=\varphi=\mathrm{id}\),
and \(J=\mathbf1=k\), \(\alpha=\mathrm{id}\).  The interval map
\(X\to J\otimes X\) is the identity, all positive derived functors
of \(\varphi\) vanish, and \(L\varphi(k)=k\).  Let
\(\Lambda(X)=X[0]\) and \(G(C)=C_0\).  These are exact adjoint
functors, since a chain map from \(X[0]\) into a nonnegative complex
is exactly a module map \(X\to C_0\).  They satisfy (134.2).
The class \(W\) consists precisely of isomorphisms of modules, so
the canonical localized functor still has only complexes concentrated
in degree zero in its essential image.  The object \(k[1]\) of
\(D_{\geq0}(k)\) has nonzero first homology and cannot lie in that
image.  Thus the asserted canonical equivalence fails.  Exactly the
missing degree-zero condition fails, since
\(\varphi(J\otimes X)=X\) need not be zero. \(\square\)

\paragraph{What is settled and what remains.}
The original Lemme 1 is proved under its weaker, explicitly annotated
hypothesis.  F-SM17 proves the exact-adjunction/counit lemma without an
infinite-coproduct assumption on its right adjoint.  BR-13403 identifies
the residual acyclic-projective test, and BR-13404 discharges it over
fields and \(\mathbf Z\) under the specified exact-adjunction data.
BR-13405 constructs those data in the locally projective ringoid case.
For a general scalar ring without the finite-injective-dimension
condition, (134.6) remains the precise test for the raw adjunction.
An equivalence constructed by a different derived or model-category
adjunction would require its own comparison with (134.5).  The much
stronger assertion that every ringoid presentation automatically
satisfies the scalar projectivity hypothesis has not been made.

